\documentclass[12pt]{article}
\usepackage[utf8]{inputenc}
\usepackage[T1]{fontenc}
\usepackage{hyperref}
\title{Algonquin Language Resources for K-12 Curriculum Development}
\author{Compiled by Claude Sonnet 5.5 (claude-sonnet-5-5) for Andr\'as Kornai}
\date{}
\begin{document}
\maketitle

\section{Introduction}
Algonquin (ISO 639-3 \texttt{alq}) is an Algic language of Quebec and Ontario, Canada, with several community dialects (e.g.\ Kitigan Zibi, Pikogan, Barriere Lake, Kitcisakik, Lac Simon, Long Point, Timiskaming). Search results were thin and partly dated (2018); several hits concern the related Anishinaabemowin/Ojibwe generally rather than Algonquin specifically.

\section{Primary Educational Resources}
\begin{itemize}
\item Algonquin College ``Kwey'' initiative, sharing Algonquin vocabulary and short videos: \url{https://www.algonquincollege.com/tri/initiatives-kwey/}
\item Blog post on Kwey at AC Online: \url{https://www.algonquincollege.com/college-blog/kwey-ac-online/}
\item Media Co-op listing of Algonquin resources, including the ``Algonquin Word of the Day'' posters for educators: \url{https://mediacoop.ca/story/algonquin-language-resources-anishinaabemowin/36702}
\end{itemize}

\section{Linguistic and Reference Materials}
\begin{itemize}
\item UBC wiki, Algonquin lexicon: \url{https://wiki.ubc.ca/Documentation:RelLex/Algonquin_Lexicon}
\item Media Co-op page (same topic, alternate link): \url{https://mediacoop.ca/node/118715}; it reports a seven-dialect Algonquin dictionary sold by the Algonquin Anishinabeg Nation Tribal Council (no direct URL for the dictionary was seen).
\end{itemize}

\section{Religious and Community Resources}
No Algonquin scripture or community organisation page with its own URL was found beyond the items above.

\section{Digital and Audio Resources}
The Media Co-op pages report daily audio lessons on CKWE 103.9FM (Kitigan Zibi dialect), some on SoundCloud, and a dictionary companion app in development as of 2018. No direct URLs were seen for these.

\section{Curriculum Development Context}
Broader Anishinaabemowin teaching guides, which concern Ojibwe and would need adaptation:
\begin{itemize}
\item Lakehead University libguide on Indigenous language teaching: \url{https://libguides.lakeheadu.ca/indigenouslangteaching}
\item Queen's University Anishinaabemowin courses: \url{https://www.queensu.ca/llcu/academics/anishinaabemowin}
\item MSU Native American resource guide, language resources: \url{https://nai.msu.edu/resources/native-american-resource-guide/native-education-legal-health-and-other-resources/language-resources}
\end{itemize}
Algonquin is taught as a subject in some schools, but no school programme or curriculum document was found in these searches. Further searching of community-run programmes is needed.

\end{document}